\documentclass[11pt]{article}
\usepackage[letterpaper,margin=0.80in,headheight=15pt]{geometry}
\usepackage[T1]{fontenc}
\usepackage{lmodern}
\usepackage{microtype}
\usepackage{amsmath,amssymb,mathtools,bm}
\usepackage{booktabs,longtable,array}
\usepackage{enumitem}
\usepackage{xcolor}
\usepackage{fancyhdr}
\usepackage{hyperref}
\allowdisplaybreaks[3]

\definecolor{navy}{HTML}{17365D}
\definecolor{teal}{HTML}{176B68}
\definecolor{softblue}{HTML}{EDF3F8}
\definecolor{softgold}{HTML}{FFF7E6}
\hypersetup{colorlinks=true,linkcolor=navy,citecolor=teal,urlcolor=teal,
 pdfauthor={Codex research report},
 pdftitle={Fully expanded bosonic Weyl and Log superinvariants from Gold's definitions}}
\setlength{\parindent}{0pt}
\setlength{\parskip}{5.2pt}
\setlist[itemize]{leftmargin=1.5em,itemsep=2pt,topsep=3pt}
\setlist[enumerate]{leftmargin=1.7em,itemsep=2pt,topsep=3pt}
\pagestyle{fancy}
\fancyhf{}
\lhead{Bosonic Weyl and Log invariants}
\rhead{Fully expanded Gold-field form}
\cfoot{\thepage}
\newcommand{\cL}{\mathcal L}
\newcommand{\cD}{\mathcal D}
\newcommand{\cR}{\mathcal R}
\newcommand{\eps}{\varepsilon}
\newcommand{\simeqbd}{\mathrel{\simeq}_{\partial}}
\newcommand{\Weyl}{\mathrm W}
\newcommand{\Log}{\log}
\newcommand{\std}{\mathrm S}
\newcommand{\Gold}{\mathrm G}
\newcommand{\Boxed}[1]{\begin{center}\fcolorbox{navy}{softblue}{\begin{minipage}{0.93\linewidth}#1\end{minipage}}\end{center}}
\newcommand{\Caution}[1]{\begin{center}\fcolorbox{teal}{softgold}{\begin{minipage}{0.93\linewidth}#1\end{minipage}}\end{center}}

\begin{document}
\pagenumbering{roman}

\begin{titlepage}
\centering
\vspace*{0.9cm}
{\color{navy}\rule{\textwidth}{1.2pt}}\\[0.60cm]
{\Huge\bfseries Bosonic Weyl and Log Superinvariants\\[0.22cm]
from Gold's Definitions}\par
\vspace{0.45cm}
{\Large Fully expanded in the original gauged dilaton--Weyl fields}\par
\vspace{0.55cm}
{\color{navy}\rule{\textwidth}{1.2pt}}\\[0.8cm]

\begin{minipage}{0.91\textwidth}
\large
This report constructs the two invariants from the BF action, the Weyl and
Log composite multiplets, their defining descent algebra, and Gold's
standard-to-dilaton--Weyl map.  The final densities contain only the original
bosonic fields and their explicit contractions.  In particular, no newly
defined two-form, composite scalar, divergence vector, or epsilon-current
symbol occurs in either final invariant.

The printed expanded Weyl and Log actions of Gold et al. are not used as
inputs or canonicalization targets.  Post-derivation comparison is used only
to document discrepancies; no printed target coefficient is imported.
\end{minipage}

\vfill
\begin{tabular}{rl}
Sector: & fermion number zero,\\
Status: & off shell and exact in $g$,\\
Gauge: & $\sigma=1$, $b_\mu=0$, compensator gaugino $=0$,\\
Representative: & boundary-equivalent; primitive currents listed below,\\
Revision: & raw Log Riemann--$w^2$ canonicalization corrected,\\
Date: & 20 August 2026.
\end{tabular}
\vspace{0.65cm}
\end{titlepage}
\setcounter{page}{2}

\begingroup
\small
\setlength{\parskip}{0pt}
\tableofcontents
\endgroup
\newpage
\pagenumbering{arabic}

\section{Problem and Target}

\subsection{Quantity being constructed}

The target is the complete fermion-number-zero part of the Weyl and Log
four-derivative superinvariants in the gauged dilaton--Weyl multiplet.  The
retained Gold fields are
\[
e_\mu{}^a,\ C_a,\ B_{ab},\ E_{abc},\ G_{ab},\ H_{ab},\
N,\ L^{ij},\ V_a{}^{ij},
\]
where $G=\mathrm dC$ and $H_{ab}$ is the dual of the gauged three-form field
strength.  All bosonic auxiliaries remain independent.  No auxiliary or
physical equation of motion, background equation, dualization, or expansion
in the gauging $g$ is permitted.

The final Weyl and Log densities are Eqs.~\eqref{eq:W-final-expanded} and
\eqref{eq:L-final-expanded}.  Every occurrence of $G$ and $H$ is displayed
with its indices or as an explicitly written scalar contraction.  There is
no definition of an auxiliary two-form such as $G-2H$, and the composite
standard-Weyl scalar $D$ has been substituted completely.

\subsection{Source provenance and quarantine}

\begin{center}
\begin{tabular}{>{\raggedright\arraybackslash}p{0.22\linewidth}>{\raggedright\arraybackslash}p{0.69\linewidth}}
\toprule
Status & Material\\
\midrule
Used & Gold et al. (2025): the BF principle, standard-to-gauged-dilaton--Weyl
map, vector identification, composite-primary definitions, three-form
definition, and gauge choice.\\
Used & Gold et al. (2023, 2024): conformal-superspace algebra, component
projections, descendant towers, degauging identities, and convention tables
needed to evaluate the defining primaries.\\
Used & The exact gauged three-form Bianchi identity derived directly in
Sec.~\ref{sec:Bianchi}.\\
Quarantined & The printed expanded Weyl and Log densities in Gold et al.
(2025), all Riemann/Gauss--Bonnet linear-combination identities, and all
scripts or notes containing those target coefficients.  They may be inspected
only after a coefficient has been fixed from the descendants.\\
Quarantined & Bosonic field equations, auxiliary elimination, background
specialization, and on-shell field redefinitions.\\
\bottomrule
\end{tabular}
\end{center}

Gold's 2025 review prints the primaries but not the complete component
descendant algebra needed by the BF density.  The 2024 component companion is
therefore indispensable definition-level input.  It supplies descendants,
not a target action used for reverse engineering.

\section{Convention Ledger}

\subsection{Geometry and indices}

Tangent indices are $a,b,\ldots=0,\ldots,4$ and $i,j=1,2$ are $SU(2)$
indices.  The tangent metric is mostly plus.  Antisymmetrization has unit
weight.  The curvature convention is
\begin{equation}
[\nabla_a,\nabla_b]X_c=R_{ab c}{}^{d}X_d,
\qquad R_{ab}=R_{ca}{}^{c}{}_{b}.
\label{eq:Rconv}
\end{equation}
The orientation is fixed by
\begin{equation}
H^{ab}=-\frac1{12}\eps^{abcde}\mathcal H_{cde}.
\label{eq:Hdual}
\end{equation}
Scalar contractions used only in intermediate arithmetic are
\[
G^2=G_{ab}G^{ab},\qquad H^2=H_{ab}H^{ab},\qquad
G\!\cdot H=G_{ab}H^{ab},\qquad L^2=L^{ij}L_{ij}.
\]
The final two equations spell these contractions out explicitly.

\subsection{Standard-component to Gold-component bridge}

Standard-superspace component fields carry a superscript $\std$; final Gold
fields carry no superscript.  The bosonic bridge is
\begin{equation}
\begin{aligned}
W^{\std}&=\sigma,&v_a^{\std}&=-C_a,&F_{ab}^{\std}&=-G_{ab},\\
X_{\std}^{ij}&=2Y^{ij}=-gL^{ij},&
w_{ab}^{\std}&=\frac{16}{3}T_{ab},&D_{\std}&=D,\\
\Phi_{ab}^{ij\,\std}&=-V_{ab}^{ij},&
\eps_{\std}^{abcde}&=-\eps^{abcde},&
\cR_{abcd}^{\std}&=-R_{abcd}.
\end{aligned}
\label{eq:bridge}
\end{equation}
For a linear multiplet,
\begin{equation}
L_{\Gold}^{ij}=-2h_{\std}^{ij},\qquad
N_{\Gold}=-F_{\std},\qquad
E_{\Gold}^a=\mathcal H_{\std}^a.
\label{eq:linearbridge}
\end{equation}
The sign $v^{\std}=-C$ follows independently by mapping the two vector
supersymmetry transformations with
$\epsilon_{\Gold}=2\xi_{\std}$,
$\Gamma_{\std}=i\gamma_{\Gold}$, and
$\lambda_{\std}=-\psi_{\Gold}$.

\subsection{Equality status and dimensions}

The notation $\cL_1\simeqbd\cL_2$ means equality after adding an explicitly
tracked covariant divergence.  Equality of actions requires a compact
manifold without boundary or boundary data for which the stated current has
vanishing normal flux.

With a four-derivative density assigned dimension four,
\begin{equation}
[\nabla]=[g]=[N]=[G]=[H]=[E]=1,
\qquad [R]=[V_{ab}]=2,\qquad [L]=[C]=0.
\label{eq:dimensions}
\end{equation}

\section{Definitions and Required Identities}

\subsection{BF principle and the two primaries}

For a vector multiplet $(\rho,A_a,Y^{ij})$ and a linear multiplet
$(L^{ij},N,E^a)$, Gold's bosonic BF density is
\begin{equation}
e^{-1}\cL_{\rm BF}=A_aE^a+\rho N+Y_{ij}L^{ij},
\qquad \nabla_aE^a=0.
\label{eq:BF-Gold}
\end{equation}
Using Eq.~\eqref{eq:linearbridge}, its standard-frame form is
\begin{equation}
e^{-1}\cL_{\rm BF}^{\std}\big|_{\rm bos}
=-\left(v_a\mathcal H^a+WF+X_{ij}h^{ij}\right)_{\std}.
\label{eq:BF-standard}
\end{equation}
This fixes the overall normalization and relative signs.

The Weyl and Log primaries are
\begin{align}
H_{\Weyl}^{ij}&=-\frac{i}{2}W^{\alpha\beta\gamma i}
W_{\alpha\beta\gamma}{}^j
+\frac{3i}{2}W^{\alpha\beta}X_{\alpha\beta}^{ij}
-\frac{3i}{4}X^{\alpha i}X_{\alpha}{}^j,
\label{eq:W-primary}\\
H_{\Log}^{ij}&=\frac{3i}{1280}
\nabla^{(ij}\nabla^{kl)}\nabla_{kl}\log W.
\label{eq:L-primary}
\end{align}
They obey the defining linearity constraint
$\nabla_\alpha^{(i}H^{jk)}=0$.

\subsection{The bosonic map, without a new two-form}

Gold's gauged three-form and standard-Weyl composites are
\begin{align}
\mathcal H_{abc}&=3\partial_{[a}B_{bc]}
+\frac32C_{[a}G_{bc]}+\frac g2E_{abc},
\label{eq:H3def}\\
T_{ab}&=\frac18G_{ab}+\frac1{48}\eps_{abcde}\mathcal H^{cde},
\label{eq:Tmap0}\\
D&=-\frac{R}{32}-\frac{G_{ab}G^{ab}}{16}
-\left(\frac{26}{3}T^{ab}-2G^{ab}\right)T_{ab}
+\frac g4N+\frac{g^2}{16}L^{ij}L_{ij}.
\label{eq:Dmap0}
\end{align}
Equation~\eqref{eq:Hdual} is substituted immediately:
\begin{equation}
T_{ab}=\frac18G_{ab}-\frac14H_{ab},
\qquad
w_{ab}^{\std}=\frac23G_{ab}-\frac43H_{ab}.
\label{eq:wmap-expanded}
\end{equation}
The scalar map follows term by term from
\begin{align}
T_{ab}T^{ab}&=\frac1{64}G_{ab}G^{ab}
-\frac1{16}G_{ab}H^{ab}+\frac1{16}H_{ab}H^{ab},\\
G^{ab}T_{ab}&=\frac18G_{ab}G^{ab}-\frac14G_{ab}H^{ab}.
\end{align}
Consequently,
\begin{equation}
\boxed{
D=-\frac{R}{32}+\frac5{96}G_{ab}G^{ab}
+\frac1{24}G_{ab}H^{ab}-\frac{13}{24}H_{ab}H^{ab}
+\frac g4N+\frac{g^2}{16}L^{ij}L_{ij}.}
\label{eq:Dmap-expanded}
\end{equation}
The right-hand side, rather than a newly named scalar, is used in every final
coefficient.

\subsection{Exact gauged three-form Bianchi identity}
\label{sec:Bianchi}

Starting with Eq.~\eqref{eq:Hdual},
\begin{equation}
\nabla_bH^{ab}=-\frac1{12}\eps^{abcde}\nabla_b\mathcal H_{cde}.
\end{equation}
The $\mathrm dB$ part vanishes by $\mathrm d^2=0$.  The three cyclic terms
in $C_{[c}G_{de]}$ are equal under the epsilon tensor, and $\mathrm dG=0$;
therefore
\begin{align}
-\frac1{12}\eps^{abcde}\nabla_b
\left(\frac32C_{[c}G_{de]}\right)
&=-\frac18\eps^{abcde}\nabla_b(C_cG_{de})\\
&=-\frac1{16}\eps^{abcde}G_{bc}G_{de}.
\end{align}
With the original constrained vector
\begin{equation}
E^a=-\frac1{12}\eps^{abcde}\nabla_bE_{cde},
\end{equation}
the exact off-shell identity is
\begin{equation}
\boxed{
\nabla_bH^{ab}=-\frac1{16}\eps^{abcde}G_{bc}G_{de}
+\frac g2E^a.}
\label{eq:Bianchi-expanded}
\end{equation}

\section{Detailed Derivation}

\subsection{Weyl invariant}

\subsubsection{Descendants and BF assembly}

At fermion number zero, descending Eq.~\eqref{eq:W-primary} gives
\begin{equation}
h_{\Weyl}^{ij}=-w^{ab}\Phi_{ab}^{ij},
\label{eq:W-lowest}
\end{equation}
and
\begin{align}
F_{\Weyl}={}&\frac{128}{3}D^2+8w^2D
-\frac23\Phi_{ab}^{ij}\Phi^{ab}_{ij}
+\frac14C_{abcd}C^{abcd}
+\frac{15}{8}C_{abcd}w^{ab}w^{cd}
\nonumber\\
&+\frac{81}{32}w_a{}^bw_b{}^cw_c{}^dw_d{}^a
-3w_{cd}\nabla_b\nabla^cw^{bd}
-\frac{33}{128}(w^2)^2
\nonumber\\
&+\frac32(\nabla^aw^{ce})(\nabla_cw_{ae})
-\frac32(\nabla^aw^{ce})(\nabla_aw_{ce})
+\frac9{16}\eps^{abcde}w_{ab}w_{cd}\nabla^fw_{ef}.
\label{eq:W-Fdesc}
\end{align}
The constrained-vector descendant is
\begin{align}
\mathcal H_{\Weyl}^a={}&\frac18\eps^{abcde}C_{bc}{}^{fg}C_{de fg}
-\frac16\eps^{abcde}\Phi_{bc}^{ij}\Phi_{de\,ij}
+\nabla_b\widetilde{\mathcal H}_{\Weyl}^{ab},
\label{eq:W-Hdesc}\\
\widetilde{\mathcal H}_{\Weyl}^{ab}={}&
8w^{ab}D+\frac32C^{abcd}w_{cd}+\frac3{16}w^{ab}w^2
-\frac94w^{ac}w_{cd}w^{db}
\nonumber\\
&-\frac94\eps^{abcde}w_{cf}\nabla_dw_e{}^f
-\frac32\eps^{abcde}w_{cf}\nabla^fw_{de}.
\label{eq:W-Htilde}
\end{align}
The divergence term is integrated once:
\begin{equation}
-v_a\nabla_b\widetilde{\mathcal H}^{ab}
=-\nabla_b(v_a\widetilde{\mathcal H}^{ab})
-\frac12F_{ab}\widetilde{\mathcal H}^{ab}.
\label{eq:BF-ibp}
\end{equation}
No field equation has entered.

\subsubsection{Algebraic expansion}

The two scalar binomials needed in the map are
\begin{align}
G^{ab}(G_{ab}-2H_{ab})&=G_{ab}G^{ab}-2G_{ab}H^{ab},\\
(G_{ab}-2H_{ab})(G^{ab}-2H^{ab})
&=G_{ab}G^{ab}-4G_{ab}H^{ab}+4H_{ab}H^{ab}.
\end{align}
For the ordered cyclic contractions, direct multiplication gives
\begin{align}
\operatorname{tr}(G-2H)^4={}&\operatorname{tr}G^4
-8\operatorname{tr}(G^3H)+16\operatorname{tr}(G^2H^2)
+8\operatorname{tr}(GHGH)
\nonumber\\
&-32\operatorname{tr}(H^3G)+16\operatorname{tr}H^4,\\
\operatorname{tr}\big((G-2H)^3G\big)={}&\operatorname{tr}G^4
-6\operatorname{tr}(G^3H)+8\operatorname{tr}(G^2H^2)
\nonumber\\
&+4\operatorname{tr}(GHGH)-8\operatorname{tr}(H^3G).
\end{align}
Here each trace is merely an intermediate matrix mnemonic.  Every chain is
written with indices in Eq.~\eqref{eq:W-final-expanded}.

The complete scalar-coefficient addition is illustrated by
\begin{align}
\cL_{\Weyl}^{R\text{-scalar}}={}&
-\frac1{12}R^2+\frac13R\,G_{ab}H^{ab}
-\frac43R\,H_{ab}H^{ab}
+\frac23gRN+\frac16g^2R L^{ij}L_{ij},
\label{eq:W-Rscalar}\\
\cL_{\Weyl}^{\text{scalar }4}={}&
-\frac1{12}(G_{ab}G^{ab})^2
-\frac13(G_{ab}G^{ab})(G_{cd}H^{cd})
+\frac83(G_{ab}G^{ab})(H_{cd}H^{cd})
\nonumber\\
&+\frac43(G_{ab}H^{ab})^2
-\frac{16}{3}(G_{ab}H^{ab})(H_{cd}H^{cd})
-4(H_{ab}H^{ab})^2,
\label{eq:W-scalar4-raw}
\end{align}
before the epsilon-square terms generated by the Bianchi reduction are
included.

\subsubsection{Canonical derivative and Bianchi reduction}

The degauged second derivatives in Eq.~\eqref{eq:W-Fdesc} are commuted only
after their displayed order is retained.  Covariant integrations by parts
and $\nabla_{[a}G_{bc]}=0$ reduce the even and parity-odd sectors.  Applying
Eq.~\eqref{eq:Bianchi-expanded} generates
\begin{align}
\cL_{\Weyl}\big|_{E}={}&-2g^2E_aE^a
+\frac{4g}{3}E_a\nabla_bG^{ab}
+\frac{5g}{12}\eps^{abcde}E_aG_{bc}G_{de}
\nonumber\\
&-\frac{2g}{3}\eps^{abcde}E_aH_{bc}H_{de}.
\label{eq:W-E-expanded}
\end{align}
Two pointwise epsilon contractions account for the associated quartic shift:
\begin{align}
&(\eps^{abcde}G_{bc}G_{de})
(\eps_{afghi}G^{fg}G^{hi})
\nonumber\\[-1mm]
&\hspace{20mm}=-8(G_{ab}G^{ab})^2
+16G_{ab}G^{bc}G_{cd}G^{da},
\label{eq:epsGGGG}\\
&(\eps^{abcde}H_{bc}H_{de})
(\eps_{afghi}G^{fg}G^{hi})
\nonumber\\[-1mm]
&\hspace{20mm}=-8(G_{ab}H^{ab})^2
+16H_{ab}H_{cd}G^{ac}G^{bd}.
\label{eq:epsHHGG}
\end{align}
Equations~\eqref{eq:W-scalar4-raw}--\eqref{eq:epsHHGG}, together with the
curvature commutators, give the one-piece Weyl density displayed in
Sec.~\ref{sec:FinalResult}.

\subsection{Log invariant}

\subsubsection{Lowest component and the corrected scalar descendant}

The bosonic lowest component of Eq.~\eqref{eq:L-primary} at $W=1$ is
\begin{align}
h_{\Log}^{ij}={}&\frac12(w^{ab}+F^{ab})\Phi_{ab}^{ij}
-\frac12\cD^a\cD_aX^{ij}
\nonumber\\
&+X^{ij}\left(-6D-\frac1{16}\cR
-\frac9{64}w^2+\frac18F^2-\frac18X^2\right).
\label{eq:L-lowest}
\end{align}
The unprojected scalar descendant contains
\begin{equation}
F_{\Log}\big|_{Y,w}=\frac3{2048}Y^2
+\frac{69}{1024}w^{ab}w_{ab}Y+\cdots,
\qquad D=-\frac3{128}Y\big|.
\label{eq:L-Yraw}
\end{equation}
Therefore $Y|=-128D/3$ and
\begin{equation}
\frac3{2048}Y^2=\frac83D^2,
\qquad
\frac{69}{1024}w^2Y=-\frac{23}{8}w^2D.
\label{eq:L-YD}
\end{equation}
The main-text value $+368w^2D/3$ applies the $Y\to D$ conversion a second
time and is not used.  The BF minus sign then contributes
$+23w^2D/8$ to the action.  This correction is fixed before, and without,
any expanded Log target density.

\subsubsection{BF assembly and raw Riemann--two-form audit}

The standard-frame BF assembly is
\begin{equation}
e^{-1}\cL_{\Log}^{\std}\big|_{\rm bos}
=-\left(v_a\mathcal H_{\Log}^a+WF_{\Log}+X_{ij}h_{\Log}^{ij}\right).
\end{equation}

The coefficient of \(\cR_{abcd}w^{ab}w^{cd}\) must be fixed before using a
canonicalized action.  We therefore start from the unsimplified bosonic
descendant/BF output of the 2024 component supplement, Sec.~3.1.  Define
\begin{equation*}
I_1:=\cR_{abcd}w^{ab}w^{cd},\qquad
I_2:=\cR_{acbd}w^{ab}w^{cd}.
\tag{L1}\label{eq:L-I12}
\end{equation*}
Contracting the first Bianchi identity
\(\cR_{abcd}+\cR_{acdb}+\cR_{adbc}=0\) with
\(w^{ab}w^{cd}\) gives
\begin{align*}
0&=I_1-I_2-I_2,
\nonumber\\
I_2&=\frac12 I_1.
\tag{L2}\label{eq:L-I2identity}
\end{align*}
In the first line, the second Bianchi term uses antisymmetry in its last
index pair, while the third uses the dummy relabeling \(c\leftrightarrow d\)
followed by \(w^{dc}=-w^{cd}\).
Let \(K:=\cR_{ab}w^{ac}w^b{}_c\).  The raw curvature block is
\begin{align*}
\cL_{w,\rm curv}^{\rm raw}
&=\frac13I_1+\frac13I_2-\frac23K-\frac13K
-\frac{27}{256}\cR\,w^{ab}w_{ab}
\nonumber\\
&=\frac12I_1-K-\frac{27}{256}\cR\,w^{ab}w_{ab}.
\tag{L3}\label{eq:L-explicit-Rww}
\end{align*}

For the derivative block define
\begin{equation*}
\mathcal N:=(\cD_aw_{bc})(\cD^aw^{bc}),\qquad
P:=(\cD^cw^a{}_b)(\cD^bw_{ac}),\qquad
Q:=(\cD^bw^a{}_b)(\cD^cw_{ac}).
\tag{L4}\label{eq:L-PQ}
\end{equation*}
Integrating the raw second derivatives once gives
\begin{align*}
\cL_{w,\rm der}^{\rm raw}
&\simeqbd
\left(\frac{39}{16}-\frac{33}{16}\right)\mathcal N
+\left(-\frac12+\frac34\right)P
+\left(2-\frac94\right)Q
\nonumber\\
&=\frac38\mathcal N+\frac14P-\frac14Q.
\tag{L5}\label{eq:L-raw-derivative-collection}
\end{align*}
Integrating one derivative at a time and using the source commutator
\(
[\cD_a,\cD_b]=-\cR_{ab}{}^{cd}M_{cd}/2+\cdots
\)
gives, with no field equation,
\begin{align*}
P&\simeqbd-w^a{}_b\cD^c\cD^bw_{ac}
\nonumber\\
&=-w^a{}_b\cD^b\cD^cw_{ac}
-w^a{}_b[\cD^c,\cD^b]w_{ac}
\nonumber\\
&\simeqbd
Q+\cR_{ab}w^{ac}w^b{}_c-\frac12I_1.
\tag{L6}\label{eq:L-PQcomm}
\end{align*}
Hence
\begin{align*}
\frac14(P-Q)
&\simeqbd
+\frac14\cR_{ab}w^{ac}w^b{}_c-\frac18I_1,
\tag{L7}\label{eq:L-PQshift}
\end{align*}
Adding Eqs.~\eqref{eq:L-explicit-Rww},
\eqref{eq:L-raw-derivative-collection}, and \eqref{eq:L-PQshift}, and then
using \(\mathcal N\simeqbd-w_{ab}\cD^2w^{ab}\), gives
\begin{equation*}
\cL_{\Log}^{\std}\big|_{w,\rm curv+kin}
=\frac38\cR_{abcd}w^{ab}w^{cd}
-\frac34\cR_{ab}w^{ac}w^b{}_c
-\frac{27}{256}\cR\,w^{ab}w_{ab}
-\frac38w^{ab}\cD^2w_{ab}.
\tag{L8}\label{eq:L-w-corrected-sector}
\end{equation*}
The same raw reduction leaves the independent \(F^{ab}F^{cd}\) coefficient
equal to \(1/4\):
\begin{equation*}
\frac13\cR_{abcd}F^{ab}F^{cd}
-\frac16\cR_{acbd}F^{ab}F^{cd}
=\left(\frac13-\frac1{12}\right)
\cR_{abcd}F^{ab}F^{cd}
=\frac14\cR_{abcd}F^{ab}F^{cd}.
\tag{L9}\label{eq:L-RFF-quarter}
\end{equation*}

Thus the complete canonicalized bosonic ledger, before the Gold map, is
\begin{align}
e^{-1}\cL_{\Log}^{\std}={}&
-\frac1{12}\eps_{\std}^{abcde}v_a\Phi_{bc}^{ij}\Phi_{de\,ij}
+\frac14\cR_{abcd}F^{ab}F^{cd}
+\frac38\cR_{abcd}w^{ab}w^{cd}
\nonumber\\
&+\frac14\cR_{ab}(4w^{ac}F^b{}_c-3w^{ac}w^b{}_c)
-\frac1{256}\cR(24F^2+80w\!\cdot F+27w^2+128D-8X^2)
\nonumber\\
&-\frac16\cR_{ab}\cR^{ab}+\frac{23}{384}\cR^2-\frac83D^2
-\frac12X^{ij}(w^{ab}\Phi_{ab\,ij}+F^{ab}\Phi_{ab\,ij})
\nonumber\\
&+\frac1{128}X^2(9w^2-8F^2+4X^2)
+\frac18D(16w\!\cdot F-8F^2+23w^2+24X^2)
\nonumber\\
&+w^{ab}F^{cd}\left(\frac3{64}w_{ab}w_{cd}
+\frac94w_{ac}w_{bd}-\frac{75}{128}w_{ab}F_{cd}
+\frac94w_{ac}F_{bd}\right.
\nonumber\\[-2mm]
&\hspace{37mm}\left.-\frac38F_{ab}F_{cd}
+\frac34F_{ac}F_{bd}\right)
\nonumber\\
&-\frac1{32}F^{ab}F^{cd}(F_{ab}F_{cd}-4F_{ac}F_{bd})
+\frac{1005}{2048}(w^2)^2+\frac16\Phi_{ab}^{ij}\Phi^{ab}_{ij}
\nonumber\\
&-\frac3{16}\eps_{\std}^{abcde}(\nabla^fw_{af})
(4w_{bc}F_{de}+F_{bc}F_{de}+3w_{bc}w_{de})
\nonumber\\
&-\frac1{16}\eps_{\std}^{abcde}(\nabla^fF_{af})
(6w_{bc}F_{de}+F_{bc}F_{de}+3w_{bc}w_{de})
\nonumber\\
&-\frac38w^{cd}\cD^a\cD_aw_{cd}
-\frac32F^{ac}\cD_c\cD^dw_{ad}
-\frac14(\cD_aX^{ij})(\cD^aX_{ij})
+\frac14(\cD_aF_{bc})(\cD^aF^{bc}).
\label{eq:L-standard-ledger}
\end{align}
This ledger is generated by the primary and derivative algebra, not read from
an expanded dilaton--Weyl invariant.

\Caution{The common coefficient \(1/4\) multiplying both \(F^{ab}F^{cd}\)
and \(w^{ab}w^{cd}\) in the printed main-text canonical form is inconsistent
with the raw descendant reduction above.  The corrected coefficients are
\(1/4\) and \(3/8\), respectively.}

\subsubsection{Direct expansion and canonicalization}

First map the corrected standard-frame Riemann terms without introducing a
new two-form.  Equations~\eqref{eq:bridge} and
\eqref{eq:wmap-expanded} give
\begin{align*}
&\frac14\cR_{abcd}^{\std}F_{\std}^{ab}F_{\std}^{cd}
+\frac38\cR_{abcd}^{\std}w_{\std}^{ab}w_{\std}^{cd}
\nonumber\\
&\quad=-\frac14R_{abcd}G^{ab}G^{cd}
-\frac16R_{abcd}(G^{ab}-2H^{ab})(G^{cd}-2H^{cd})
\nonumber\\
&\quad=-\frac5{12}R_{abcd}G^{ab}G^{cd}
+\frac23R_{abcd}G^{ab}H^{cd}
-\frac23R_{abcd}H^{ab}H^{cd}.
\tag{L10}\label{eq:L-corrected-explicit-Riemann}
\end{align*}
The two mixed products combine because
\(R_{abcd}=R_{cdab}\).  Applying the commutator identity
\eqref{eq:twoform-boundary-id} to the remaining mapped derivative terms gives
the curvature shift
\begin{equation*}
+\frac5{12}R_{abcd}G^{ab}G^{cd}
-\frac23R_{abcd}G^{ab}H^{cd}
+\frac23R_{abcd}H^{ab}H^{cd}.
\tag{L11}\label{eq:L-derivative-Riemann-shift}
\end{equation*}
The complete definition-led Riemann--two-form ledger is consequently
\begin{equation*}
\begin{array}{c|rrr}
&R_{abcd}G^{ab}G^{cd}&R_{abcd}G^{ab}H^{cd}
&R_{abcd}H^{ab}H^{cd}\\
\midrule
\text{explicit corrected Riemann terms}
&-\frac5{12}&\frac23&-\frac23\\
\text{derivative commutator shift}
&\frac5{12}&-\frac23&\frac23\\
\midrule
\text{corrected total}&0&0&0
\end{array}
\tag{L12}\label{eq:L-Riemann-cancellation-ledger}
\end{equation*}
For comparison only after this derivation, the change relative to the earlier
\(1/4\)-\(1/4\) intermediate ledger is
\begin{align*}
\delta\cL_{\Log}
&=\frac18\cR_{abcd}^{\std}w_{\std}^{ab}w_{\std}^{cd}
\nonumber\\
&=-\frac1{18}R_{abcd}G^{ab}G^{cd}
+\frac29R_{abcd}G^{ab}H^{cd}
-\frac29R_{abcd}H^{ab}H^{cd}.
\tag{L13}\label{eq:L-Riemann-correction-map}
\end{align*}
This is a pointwise algebraic correction; it introduces no additional
boundary current.

Substituting Eqs.~\eqref{eq:bridge}, \eqref{eq:wmap-expanded}, and
\eqref{eq:Dmap-expanded} first gives the derivative representative
\begin{align}
\cL_{\Log}^{\nabla,\rm raw}={}&
\frac5{12}(\nabla_aG_{bc})(\nabla^aG^{bc})
-\frac23(\nabla_aG_{bc})(\nabla^aH^{bc})
+\frac23(\nabla_aH_{bc})(\nabla^aH^{bc})
\nonumber\\
&-(\nabla_bG^{ab})(\nabla^cG_{ac})
+2(\nabla_bG^{ab})(\nabla^cH_{ac})
-\frac{g^2}{4}(\nabla_aL^{ij})(\nabla^aL_{ij}),
\label{eq:L-der-raw}
\end{align}
and the parity-odd representative
\begin{align}
\cL_{\Log}^{\rm odd,raw}={}&
+\frac1{16}(\nabla^bG_{ab})\eps^{acdef}G_{cd}G_{ef}
+\frac1{12}(\nabla^bH_{ab})\eps^{acdef}G_{cd}G_{ef}
\nonumber\\
&+\frac13(\nabla^bG_{ab})\eps^{acdef}H_{cd}H_{ef}
-\frac43(\nabla^bH_{ab})\eps^{acdef}H_{cd}H_{ef}
\nonumber\\
&-\frac16(\nabla^bG_{ab})\eps^{acdef}H_{cd}G_{ef}.
\label{eq:L-odd-raw}
\end{align}
Using the two-form commutator identity, $\nabla_{[a}G_{bc]}=0$, and then
Eq.~\eqref{eq:Bianchi-expanded}, Eqs.~\eqref{eq:L-der-raw} and
\eqref{eq:L-odd-raw} become the derivative, curvature-shift, double-epsilon,
and explicit $E_a$ terms in Eq.~\eqref{eq:L-final-expanded}.

For transparency, the scalar quartic coefficients arise from the exact sums
\begin{equation}
\begin{array}{c|rrrrrr}
\text{monomial}&(G^2)^2&G^2(G\cdot H)&G^2H^2&(G\cdot H)^2&(G\cdot H)H^2&(H^2)^2\\
\midrule
\text{coefficient}&-\frac1{48}&-\frac13&\frac23&\frac43&-\frac43&-2
\end{array}
\label{eq:L-scalar-ledger}
\end{equation}
and the ordered cyclic-chain coefficients are
\begin{equation}
\begin{array}{c|rrrrr}
\text{chain}&GGGG&GGGH&GGHH&GHGH&GHHH\\
\midrule
\text{coefficient}&-\frac1{24}&1&-\frac43&-\frac83&\frac{16}{3}.
\end{array}
\label{eq:L-trace-ledger}
\end{equation}
All chains are displayed with indices in the final result.

\section{Consistency Checks}

\subsection{Exact-rational and indexed-tensor audit}

The independent script
\begin{center}
\texttt{scripts/verify\_weyl\_log\_definition\_map.py}
\end{center}
contains no coefficients from the quarantined final actions.  It performs:
\begin{enumerate}
\item exact reconstruction of the complete $D$ composite;
\item the unprojected Log $Y\to D$ conversion;
\item exact raw-supplement canonicalization of the quadratic Log $w$ sector,
including the $1/4\to3/8$ Riemann coefficient and its fully mapped
cancellation;
\item complete polynomial equality, including absence of unlisted monomials,
for both scalar sectors;
\item exact indexed five-dimensional tests of all noncommuting quartic chains;
\item exact flat/local derivative-order and epsilon-monomial expansions;
\item the two epsilon-square identities;
\item the full gauged Bianchi substitutions and both $E_3$ sectors;
\item mass-dimension and limiting-case checks; and
\item a lexical gate forbidding newly defined shorthand in the two final
densities.
\end{enumerate}

\subsection{Dimensional, index, and limiting checks}

Every free index in Eqs.~\eqref{eq:W-final-expanded} and
\eqref{eq:L-final-expanded} is saturated.  Using Eq.~\eqref{eq:dimensions},
representative dimensions are
\[
[R_{abcd}R^{abcd}]=4,\quad[gRN]=4,\quad
[g^2(\nabla L)^2]=4,\quad[g^2E^2]=4,
\quad[g\eps EGG]=4.
\]
The following substitutions were tested directly:
\begin{enumerate}
\item $G=H=N=L=V=E=0$ gives
$\cL_{\Weyl^2}=-R_{abcd}^2/4+R_{ab}^2/3-R^2/12$ and
$\cL_{\Log}=-R_{ab}^2/6+R^2/24$;
\item $g=0$ removes every explicit $E_3$ term while leaving the off-shell
$G,H$ sectors;
\item for the correction in Eq.~\eqref{eq:L-Riemann-correction-map}, $H=0$ gives
$-R_{abcd}G^{ab}G^{cd}/18$, $G=0$ gives
$-2R_{abcd}H^{ab}H^{cd}/9$, and $G=2H$ gives zero; adding these to the
earlier intermediate block gives zero identically;
\item flat space with constant forms removes every derivative and curvature
monomial; and
\item setting all fields to zero makes both densities vanish.
\end{enumerate}

\subsection{Boundary ledger}

The Weyl density uses the BF integration by parts
Eq.~\eqref{eq:BF-ibp}, derivative canonicalization, and the explicitly
expanded current in Appendix~\ref{app:Wcurrent}.  The Log density additionally
uses $-w\Box w\simeqbd(\nabla w)^2$ and the closed-$G$ two-form commutator.
The primitive Log currents are as follows.

First, if the constrained-vector descendant is written as
$\mathcal H_{\Log}^a=\mathcal H_{\Log,0}^a+\nabla_bQ_{\Log}^{ab}$ with
$Q_{\Log}^{ab}=-Q_{\Log}^{ba}$, the BF integration is
\begin{equation}
-v_a\nabla_bQ_{\Log}^{ab}
=-\nabla_b(v_aQ_{\Log}^{ab})-\frac12F_{ab}Q_{\Log}^{ab},
\qquad \mathcal J_{\rm BF}^b=-v_aQ_{\Log}^{ab}.
\end{equation}
Here $Q_{\Log}^{ab}$ is the defining antisymmetric descendant potential, not
a new physical field.

Second, the mapped $-w\Box w$ operation has current
\begin{align}
\mathcal J_w^a={}&-\frac16G^{cd}\nabla^aG_{cd}
+\frac13G^{cd}\nabla^aH_{cd}
+\frac13H^{cd}\nabla^aG_{cd}
-\frac23H^{cd}\nabla^aH_{cd}.
\end{align}
Third, the ordered mixed second derivative is treated by
\begin{align}
&G^{ac}\nabla_c\nabla^d(G_{ad}-2H_{ad})
\nonumber\\
&\quad=-(\nabla_cG^{ac})\nabla^d(G_{ad}-2H_{ad})
+\nabla_c\!\left[G^{ac}\nabla^d(G_{ad}-2H_{ad})\right],
\label{eq:Log-mixed-current}
\end{align}
so its current is the expression inside the square brackets.

Finally, the two-form canonicalization follows from the general identity
\begin{align}
(\nabla^aX^{ce})(\nabla_cY_{ae})\simeqbd{}&
(\nabla_bX^{ab})(\nabla^cY_{ac})
-R_{ab}X^{ac}Y^b{}_c
+\frac12R_{abcd}X^{ab}Y^{cd},
\label{eq:twoform-boundary-id}
\end{align}
whose primitive current is
\begin{equation}
\mathcal J_{X,Y}^m=(\nabla^aX^{me})Y_{ae}
+Y^m{}_e\nabla_bX^{eb}.
\end{equation}
Equation~\eqref{eq:twoform-boundary-id} is applied only to the explicitly
listed pairs $X,Y\in\{G,H\}$; $X,Y$ are placeholders in the identity, not
additional fields.

The ordinary Levi--Civita divergence descending from $-4\Box D$ is dropped as
a whole:
\begin{equation}
\nabla^2\!\left[
\frac18R-\frac5{24}G_{ab}G^{ab}-\frac16G_{ab}H^{ab}
+\frac{13}{6}H_{ab}H^{ab}-gN
-\frac{g^2}{4}L^{ij}L_{ij}\right].
\label{eq:Log-D-boundary}
\end{equation}
The corresponding primitive current is $\nabla^a$ of the square bracket in
Eq.~\eqref{eq:Log-D-boundary}.
Its conformal curvature-weight contribution is retained.  In particular, one
must not keep only the $-g\nabla^2N$ piece while discarding the rest of
Eq.~\eqref{eq:Log-D-boundary}.

\section{Final Result}
\label{sec:FinalResult}

\subsection{Fully expanded Weyl invariant}

After the full gauged Bianchi identity, curvature commutators, and the
boundary operations listed above, the one-piece Gold-field representative is
\begingroup
\small
\begin{align}
e^{-1}\cL_{\Weyl^2}^{\rm bos}\simeqbd{}&
-\frac18\eps^{abcde}C_aR_{bc}{}^{fg}R_{de fg}
+\frac16\eps^{abcde}C_aV_{bc}^{ij}V_{de\,ij}
+\frac23V_{ab}^{ij}V^{ab}_{ij}
\nonumber\\
&-\frac14R_{abcd}R^{abcd}+\frac13R_{ab}R^{ab}-\frac1{12}R^2
\nonumber\\
&+\frac13R_{abcd}G^{ab}G^{cd}
-R_{abcd}H^{ab}G^{cd}
-\frac23R_{abcd}H^{ab}H^{cd}
\nonumber\\
&-\frac43R_{ab}H^{ac}G^b{}_c
+\frac{16}{3}R_{ab}H^{ac}H^b{}_c
+\frac13R\,H_{ab}G^{ab}-\frac43R\,H_{ab}H^{ab}
\nonumber\\
&+\frac23gRN+\frac16g^2R L^{ij}L_{ij}
\nonumber\\
&+\frac1{12}(G_{ab}G^{ab})^2
-\frac13(G_{ab}G^{ab})(G_{cd}H^{cd})
+\frac83(G_{ab}G^{ab})(H_{cd}H^{cd})
\nonumber\\
&+\frac23(G_{ab}H^{ab})^2
-\frac{16}{3}(G_{ab}H^{ab})(H_{cd}H^{cd})
-4(H_{ab}H^{ab})^2
\nonumber\\
&-\frac12G_{ab}G^{bc}G_{cd}G^{da}
+2G_{ab}G^{bc}G_{cd}H^{da}
-\frac{16}{3}G_{ab}G^{bc}H_{cd}H^{da}
\nonumber\\
&-\frac43G_{ab}H^{bc}G_{cd}H^{da}
+\frac{40}{3}G_{ab}H^{bc}H_{cd}H^{da}
-8H_{ab}H^{bc}H_{cd}H^{da}
\nonumber\\
&-\frac43gN\,G_{ab}G^{ab}
+\frac43gN\,G_{ab}H^{ab}
+8gN\,H_{ab}H^{ab}
\nonumber\\
&-\frac13g^2L^{ij}L_{ij}\,G_{ab}G^{ab}
+\frac13g^2L^{ij}L_{ij}\,G_{ab}H^{ab}
+2g^2L^{ij}L_{ij}\,H_{ab}H^{ab}
\nonumber\\
&-\frac83g^2N^2-\frac43g^3N L^{ij}L_{ij}
-\frac16g^4(L^{ij}L_{ij})^2
\nonumber\\
&-\frac13(\nabla_aG_{bc})(\nabla^aG^{bc})
+\frac83(\nabla_aH_{bc})(\nabla^aH^{bc})
\nonumber\\
&+\frac43\eps^{abcde}H_{ab}H_{cd}\nabla^fG_{ef}
-2\eps^{abcde}H_b{}^f(\nabla_aH_{cf})G_{de}
\nonumber\\
&-\frac23\eps^{abcde}H_{ab}(\nabla^fG_{cf})G_{de}
-\frac1{24}\eps^{abcde}(\nabla^fG_{af})G_{bc}G_{de}
\nonumber\\
&-2g^2E_aE^a+\frac{4g}{3}E_a\nabla_bG^{ab}
+\frac{5g}{12}\eps^{abcde}E_aG_{bc}G_{de}
-\frac{2g}{3}\eps^{abcde}E_aH_{bc}H_{de}
\nonumber\\
&+\frac23gV_{ab}^{ij}L_{ij}G^{ab}
-\frac43gV_{ab}^{ij}L_{ij}H^{ab}.
\label{eq:W-final-expanded}
\end{align}
\endgroup

\subsection{Fully expanded Log invariant}

In the same original field basis and Bianchi-resolved representative,
\begingroup
\small
\begin{align}
e^{-1}\cL_{\Log}^{\rm bos}\simeqbd{}&
-\frac1{12}\eps^{abcde}C_aV_{bc}^{ij}V_{de\,ij}
+\frac16V_{ab}^{ij}V^{ab}_{ij}
+\frac g6V_{ab}^{ij}L_{ij}G^{ab}
+\frac{2g}{3}V_{ab}^{ij}L_{ij}H^{ab}
\nonumber\\
&-\frac16R_{ab}R^{ab}+\frac1{24}R^2
+\frac16R_{ab}G^{ac}G^b{}_c
-\frac43R_{ab}G^{ac}H^b{}_c
\nonumber\\
&+R\left(\frac g6N-\frac{g^2}{12}L^{ij}L_{ij}
+\frac13G_{ab}H^{ab}-\frac13H_{ab}H^{ab}\right)
\nonumber\\
&-\frac{g^2}{6}N^2+\frac{2g^3}{3}N L^{ij}L_{ij}
+\frac{5g^4}{24}(L^{ij}L_{ij})^2
\nonumber\\
&+gN\left(-\frac13G_{ab}G^{ab}-\frac23G_{ab}H^{ab}
+2H_{ab}H^{ab}\right)
\nonumber\\
&+g^2L^{ij}L_{ij}\left(\frac1{24}G_{ab}G^{ab}
-\frac16G_{ab}H^{ab}-H_{ab}H^{ab}\right)
\nonumber\\
&-\frac1{48}(G_{ab}G^{ab})^2
-\frac13(G_{ab}G^{ab})(G_{cd}H^{cd})
+\frac23(G_{ab}G^{ab})(H_{cd}H^{cd})
\nonumber\\
&+\frac43(G_{ab}H^{ab})^2
-\frac43(G_{ab}H^{ab})(H_{cd}H^{cd})
-2(H_{ab}H^{ab})^2
\nonumber\\
&-\frac1{24}G_{ab}G^{bc}G_{cd}G^{da}
+G_{ab}G^{bc}G_{cd}H^{da}
-\frac43G_{ab}G^{bc}H_{cd}H^{da}
\nonumber\\
&-\frac83G_{ab}H^{bc}G_{cd}H^{da}
+\frac{16}{3}G_{ab}H^{bc}H_{cd}H^{da}
\nonumber\\
&-\frac16(\nabla_bG^{ab})(\nabla^cG_{ac})
+\frac23(\nabla_aH_{bc})
\left(\nabla^aH^{bc}+\nabla^bH^{ca}+\nabla^cH^{ab}\right)
\nonumber\\
&-\frac{g^2}{4}(\nabla_aL^{ij})(\nabla^aL_{ij})
\nonumber\\
&+\frac1{48}(\nabla^bG_{ab})\eps^{acdef}G_{cd}G_{ef}
+\frac13(\nabla^bG_{ab})\eps^{acdef}H_{cd}H_{ef}
\nonumber\\
&-\frac16(\nabla^bG_{ab})\eps^{acdef}H_{cd}G_{ef}
\nonumber\\
&+\frac1{12}(\eps^{abcde}G_{bc}G_{de})
(\eps_{a}{}^{fghi}H_{fg}H_{hi})
\nonumber\\
&+\frac g3E_a\nabla_bG^{ab}
-\frac g{24}E_a\eps^{abcde}G_{bc}G_{de}
-\frac{2g}{3}E_a\eps^{abcde}H_{bc}H_{de}
+\frac{g^2}{3}E_aE^a.
\label{eq:L-final-expanded}
\end{align}
\endgroup

The absence of explicit
\(R_{abcd}G^{ab}G^{cd}\),
\(R_{abcd}G^{ab}H^{cd}\), and
\(R_{abcd}H^{ab}H^{cd}\) terms in
Eq.~\eqref{eq:L-final-expanded} is the exact cancellation displayed in
Eq.~\eqref{eq:L-Riemann-cancellation-ledger}, not a truncation.

\Boxed{Equations \eqref{eq:W-final-expanded} and
\eqref{eq:L-final-expanded} are the requested complete bosonic invariants.
They contain only Gold's original fields and explicit contractions.  Their
$E_3$ terms are already included through the exact identity
\eqref{eq:Bianchi-expanded}; they must not be added a second time.  In the
corrected Log density, the definition-level $3/8$ coefficient removes the
entire explicit Riemann--$G,H$ block.}

\section{Convention Dependence and Remaining Uncertainties}

\begin{enumerate}
\item \textbf{Bosonic scope.}  The result is complete only at fermion number
zero.  No all-fermion claim is made.
\item \textbf{Source closure.}  Gold's 2025 review alone does not print the
full descendants needed by the BF action.  The 2023/2024 derivative algebra
and component projections are necessary definition-level inputs.
\item \textbf{Resolved Log component issues.}  The unprojected descendant
and the defining $D=-3Y/128$ projection fix
$F_{\Log}\supset-23w^2D/8$.  Independently, canonicalizing the raw quadratic
$w$ sector fixes
$\cR_{abcd}F^{ab}F^{cd}/4+3\cR_{abcd}w^{ab}w^{cd}/8$ and gives
Eq.~\eqref{eq:L-Riemann-correction-map}.  Confidence: certain.
\item \textbf{Boundary representative.}  Both displayed densities are
boundary-equivalent to their BF representatives.  A physical boundary
requires restoration of the listed currents and a boundary-condition or
counterterm analysis.
\item \textbf{Orientation and curvature signs.}  Reversing the epsilon
orientation or Riemann convention changes the parity-odd or curvature bridge.
Equations \eqref{eq:Rconv}, \eqref{eq:Hdual}, and \eqref{eq:bridge} are part
of the result.
\item \textbf{No target fitting.}  The expanded final equations of Gold
et al. were not used as derivation inputs or canonicalization targets.
Post-derivation comparison is used only to document the two component-level
discrepancies above.  The formulas remain a direct definition-led
representative, not a coefficient fit to a published boundary basis.
\end{enumerate}

\appendix
\section{Expanded Weyl Boundary Current}
\label{app:Wcurrent}

The local BF density and Eq.~\eqref{eq:W-final-expanded} satisfy
\begin{equation}
e^{-1}\cL_{\rm local}=e^{-1}\cL_{\Weyl^2}^{\rm bos}
+\nabla_m\mathcal J^m,
\qquad
\mathcal J^m=C_a\Theta^{am}+\mathcal J_{\rm even}^m
+\mathcal J_{\rm odd}^m.
\end{equation}
All pieces are written in the original fields:
\begin{align}
\Theta^{ab}={}&-R^{abcd}G_{cd}+2R^{abcd}H_{cd}
-\frac43R^{[a}{}_cG^{b]c}+\frac83R^{[a}{}_cH^{b]c}
-\frac13RG^{ab}+\frac23RH^{ab}
\nonumber\\
&+\frac13G^{ab}G_{cd}G^{cd}-\frac83G^{ab}H_{cd}H^{cd}
-\frac23H^{ab}G_{cd}G^{cd}+\frac{16}{3}H^{ab}H_{cd}H^{cd}
\nonumber\\
&+\frac43gNG^{ab}-\frac83gNH^{ab}
+\frac13g^2L^{ij}L_{ij}G^{ab}
-\frac23g^2L^{ij}L_{ij}H^{ab}
\nonumber\\
&-\frac23G^{ac}G_{cd}G^{db}
+\frac43\left(G^{ac}G_{cd}H^{db}+G^{ac}H_{cd}G^{db}
+H^{ac}G_{cd}G^{db}\right)
\nonumber\\
&-\frac83\left(G^{ac}H_{cd}H^{db}+H^{ac}G_{cd}H^{db}
+H^{ac}H_{cd}G^{db}\right)
+\frac{16}{3}H^{ac}H_{cd}H^{db}
\nonumber\\
&+\eps^{abcde}\left(G_{cf}\nabla_dG_e{}^f
-2G_{cf}\nabla_dH_e{}^f-2H_{cf}\nabla_dG_e{}^f
+4H_{cf}\nabla_dH_e{}^f\right)
\nonumber\\
&+\frac23\eps^{abcde}\left(G_{cf}\nabla^fG_{de}
-2G_{cf}\nabla^fH_{de}-2H_{cf}\nabla^fG_{de}
+4H_{cf}\nabla^fH_{de}\right),
\end{align}
\begin{align}
\mathcal J_{\rm even}^m={}&
\frac43G^m{}_d\nabla_bG^{bd}-\frac83G^m{}_d\nabla_bH^{bd}
-\frac83H^m{}_d\nabla_bG^{bd}
+\frac{16}{3}H^m{}_d\nabla_bH^{bd}
\nonumber\\
&+\frac43\left(G^{ce}\nabla_cG^m{}_e+G^{me}\nabla_bG_e{}^b\right)
\nonumber\\
&-\frac83\left(G^{ce}\nabla_cH^m{}_e+G^{me}\nabla_bH_e{}^b\right)
\nonumber\\
&-\frac83\left(H^{ce}\nabla_cH^m{}_e+H^{me}\nabla_bH_e{}^b\right),\\
\mathcal J_{\rm odd}^m={}&
\frac13\eps^{abcme}G_{ab}H_{cf}G_e{}^f
+\frac43\eps^{abcde}G_{ab}H_{cd}H_e{}^m.
\end{align}
Thus the action equality assumes $n_m\mathcal J^m=0$ at the boundary.

\section{Source Anchors}

\begin{longtable}{>{\raggedright\arraybackslash}p{0.27\linewidth}>{\raggedright\arraybackslash}p{0.65\linewidth}}
\toprule
Source & Exact role\\
\midrule
\endfirsthead
\toprule Source & Exact role\\ \midrule
\endhead
Gold et al., \emph{On five-dimensional curvature squared supergravity and
holography}, arXiv:2504.10856, Secs.~2.2--2.3 & Gauged composite map,
$\mathcal H_3$, BF principle, vector identification, primaries, and gauge.
Its printed expanded final invariant equations are quarantined.\\
Gold et al., \emph{On curvature-squared invariants of minimal
five-dimensional supergravity from superspace}, arXiv:2302.14295, Secs.~3--4
& Superspace definitions of the Weyl and Log composites and BF construction.\\
Gold et al., \emph{Components of curvature-squared invariants of minimal
supergravity in five dimensions}, arXiv:2311.00679, Secs.~2, 4 and
Apps.~A--C & Component projections, full derivative algebra, descendant
towers, degauging identities, and convention bridge.\\
Gold et al. component supplement, Secs.~2.2.2 and 3.1 & Unprojected Log
scalar descendant $3Y^2/2048+69w^2Y/1024+\cdots$; raw quadratic $w$ block
and ordered-derivative canonicalization fixing
$3\cR_{abcd}w^{ab}w^{cd}/8$.\\
Coomans--Ozkan, arXiv:1210.4704, Sec.~5 & Independent definition anchor for
the gauged dilaton--Weyl map and exact three-form Bianchi chain.\\
\bottomrule
\end{longtable}

\section{Reproduction Commands}

From the project root, the definition-map audit is reproduced with
\begin{verbatim}
python3 scripts/verify_weyl_log_definition_map.py
\end{verbatim}
The editable report source is
\begin{verbatim}
results/latex/weyl_log_from_definitions/weyl_log_from_definitions.tex
\end{verbatim}
and TeX build artifacts are stored only under
\begin{verbatim}
results/build/weyl_log_from_definitions/
\end{verbatim}

\end{document}
